\documentclass[11pt]{article}
\usepackage[margin=1in]{geometry}
\usepackage{amsmath,amssymb,amsthm,mathtools}
\usepackage[T1]{fontenc}
\usepackage[expansion=false]{microtype}
\usepackage[hidelinks]{hyperref}
\newtheorem{theorem}{Theorem}[section]
\newtheorem{lemma}[theorem]{Lemma}
\newtheorem{proposition}[theorem]{Proposition}
\newtheorem{remark}[theorem]{Remark}
\newcommand{\ddc}{\mathrm{d}\mathrm{d}^c}
\newcommand{\MA}{\operatorname{MA}}
\newcommand{\suppset}{\mathcal A}
\title{Weak adjoints for logarithmic models on non-Gauduchon product surfaces}
\author{}
\date{7 October 2026}
\begin{document}
\maketitle
\begin{abstract}
We give explicit sufficient conditions for a weak adjoint in a genuinely
singular ordinary reference type on $\mathbb P^1\times\mathbb P^1$.
The background is nonclosed and none of its smooth positive potential
perturbations is Gauduchon. A logarithmic divisor has an explicit admissible
ordinary model. For a family of continuous logarithmic densities, the
solution has continuous, possibly nonsmooth, positive residual coefficients.
Its directly defined weighted non-pluripolar current is closed. We prove
uniqueness of the supporting probability on the full space of continuous
data by a signed-measure resolvent argument. Every fixed-divisor approximation
with uniformly convergent positive residual forms has the same probability,
weighted-current and weak-weight limits, without symmetry or derivative
bounds. An explicit nonsmooth example distinguishes the support from
normalized Monge--Amp\`ere volume. Filling the pole by smooth potentials
instead retains positive divisor mass, explaining the approximation boundary.
\end{abstract}

\section{Framework and the sufficient conditions}

We use $\ddc=(i/(2\pi))\partial\bar\partial$. Potentials have genuine
canonical upper values. Write $v\preceq\phi$ when $v\le\phi+C$, and
$v\sim\phi$ for mutual bounded comparisons, interpreted as extended-value
inequalities at common poles. For the canonical obstacle envelope $P_\theta$,
put
\[
 P_\theta[\phi]=\left(\lim_{C\to\infty}
 P_\theta(\min\{\phi+C,0\})\right)^*.
\]
An ordinary model satisfies $\sup\phi=0$ and $P_\theta[\phi]=\phi$.
The admissible reference framework used here supplies bounded-density
solvability and uniqueness of the scalar, after the additional envelope
$P_{\epsilon\theta}(\phi)$ is genuine for some $0<\epsilon<1$.
Thus, for bounded logarithmic data $G$, an ordinary-type equation has scalar
$c_\phi(G)>0$, and $B_\phi(G)=\log c_\phi(G)$. For continuous data the full
supporting set is nonempty and consists of probabilities satisfying
\begin{equation}\label{eq:support}
 B_\phi(G)\le B_\phi(F)-\int_X(G-F)\,\mathrm{d}\mu
 \quad\text{for every }G\in C(X).
\end{equation}
These are the given admissible normalization framework, rather than new
singular existence or support theorems. Our task is to prove the additional
weak product, operator, uniqueness and approximation assertions. Products
below are non-pluripolar products with the smooth-background corrections.

Let $X=\mathbb P^1_x\times\mathbb P^1_y$, with normalized Fubini--Study
forms $\omega_i$, $\int_{\mathbb P^1}\omega_i=1$. Put
$\mathrm{d}W=\omega_1\wedge\omega_2$ and $\mathrm{d}V=q\,\mathrm{d}W$,
where $q$ is smooth and strictly positive and $\int_Xq\,\mathrm{d}W=1$.
Choose $a,\lambda>0$ and a smooth strictly positive nonconstant function
$b=b(x)$. Define
\[
 K=a+\lambda,\quad \theta=K\omega_1+b\omega_2,
 \quad\eta=a\omega_1+b\omega_2,
 \quad D=\{Z_0=0\}\times\mathbb P^1_y,
\]
\[
 \psi=\lambda\log\frac{|Z_0|^2}{|Z_0|^2+|Z_1|^2}.
\]
Both smooth backgrounds are positive. In particular $\theta$ is big.
Poincar\'e--Lelong gives
\begin{equation}\label{eq:pole}
 \theta+\ddc\psi=\eta+\lambda[D].
\end{equation}
Set $R=a/\lambda$ and $C_*=K\log(K/\lambda)-a\log(a/\lambda)$.
In the chart $z=Z_1/Z_0$ define
\begin{equation}\label{eq:model}
 \phi(z)=\begin{cases}
 0,&|z|^2\le R,\\
 a\log|z|^2+C_*-K\log(1+|z|^2),&|z|^2\ge R,
 \end{cases}
 \qquad\phi|_D=-\infty.
\end{equation}

\begin{theorem}\label{thm:main}
The function $\phi$ is an admissible ordinary model with $\phi\sim\psi$,
positive finite mass, and a residual horizontal form vanishing on an open
set. No smooth positive potential perturbation of $\theta$ is Gauduchon.

Let $F\in C(X)$ satisfy
\begin{equation}\label{eq:density}
 \frac{q(x,y)e^{F(x,y)}}{2b(x)}=H(x)>0,
 \qquad H\in C(\mathbb P^1_x).
\end{equation}
Put $m=\int H\omega_1$, $c=a/m$, and $A=cH$. There is a mean-zero
bounded continuous real $w=w(x)$ with $\ddc w=(A-a)\omega_1$.
The actual singular solution $u=\psi+w$ has $u\sim\phi$ and
\[
 \langle\theta+\ddc u\rangle=\Omega=A\omega_1+b\omega_2,
 \qquad\MA^{\mathrm{ext}}_\theta(u)=\Omega^2=c e^F\mathrm{d}V.
\]
In particular $c_\phi(F)=c$. The directly defined weight and current
\begin{equation}\label{eq:weight}
 \rho=\frac1{2ab},\qquad Q=\rho\Omega
 =\frac{A}{2ab}\omega_1+\frac1{2a}\omega_2
\end{equation}
satisfy $\mathrm{d}Q=0$ and $\int\rho\Omega^2=1$.
For every smooth real $h$ the weak and logarithmic linearizations are
\[
 D_u\MA^{\mathrm{ext}}(h)=2\Omega\wedge\ddc h,
 \qquad L_uh=\frac{2\Omega\wedge\ddc h}{\Omega^2}
 =\frac{\Delta_1h}{A}+\frac{\Delta_2h}{b},
\]
where $\ddc_i h=(\Delta_i h)\omega_i$. The unique invariant probability for $L_u$,
and the unique full supporting probability, are
\begin{equation}\label{eq:nu}
 \nu=\rho\Omega^2=\frac A a\,\mathrm{d}W,
 \qquad \suppset_\phi(F)=\{\nu\}.
\end{equation}
Uniqueness here applies to all Borel invariant probabilities, and, almost everywhere, to all
nonnegative integrable weights normalized by $\int r\Omega^2=1$ and
satisfying $\ddc(r\Omega)=0$.

For any smooth real $w_j$ on $X$ with
\begin{equation}\label{eq:approx}
 \Omega_j=\eta+\ddc w_j>0,
 \qquad\|\Omega_j-\Omega\|_{C^0,\omega_1+\omega_2}\longrightarrow0,
\end{equation}
retain the exact pole by setting $u_j=\psi+w_j$, and set
$F_j=\log(\Omega_j^2/(c\,\mathrm{d}V))$. Normalize the classical residual
adjoint by
\[
 \ddc(\rho_j\Omega_j)=0,\quad\rho_j>0,
 \qquad\int_X\rho_j\Omega_j^2=1.
\]
Then $F_j\to F$ uniformly and, for the entire sequence,
\begin{equation}\label{eq:limits}
 \rho_j\Omega_j^2\rightharpoonup\nu,\qquad
 \rho_j\Omega_j\rightharpoonup Q,\qquad
 \rho_j\,\mathrm{d}W\rightharpoonup\rho\,\mathrm{d}W.
\end{equation}
The approximation class is nonempty with $w_j\to w$ uniformly and permits
nonsymmetric sequences. It requires no bounds on higher derivatives or
pointwise convergence of the weights.
\end{theorem}

The theorem concerns this constructed solution in a genuine unbounded
ordinary type. It does not assert uniqueness of every rough solution
potential, degenerate solution residuals, arbitrary densities or singularity
classes, or an adjoint limit under potential convergence alone.

\section{Ordinary model and genuine admissibility}

For $C\ge0$, $\psi$ is a candidate in
$E_C=P_\theta(\min\{\psi+C,0\})$. The transitive $\mathrm{PU}(2)$ action
on the second sphere preserves the background and obstacle. Equivariance
of the canonical envelope makes $E_C$ independent of $y$. For a function
of $x$, positivity is precisely $K\omega_1+\ddc v\ge0$, since its second
block $b\omega_2$ is positive. Rotation makes the one-dimensional envelope
radial. Its total local potential
$g(s)=v+K\log(1+e^s)$, $s=\log|z|^2$, is convex and satisfies
\[
 g(s)\le K\log(1+e^s),\qquad g(s)\le a\log(1+e^s)+C.
\]
The latter bounds every one-sided slope by $a$. At $s_0=\log R$, therefore,
\[
 g(s)\le K\log(1+e^{s_0})+a(s-s_0)\quad(s\ge s_0).
\]
For $s\le s_0$ use the first obstacle. The maximal candidate follows that
obstacle up to $s_0$, then its tangent line of slope $a$. Subtracting
$K\log(1+e^s)$ gives \eqref{eq:model}. Convexity and matching first
derivatives prove positivity with no negative boundary current. Moreover
\[
 0\le\phi-\psi\le C_*\quad\text{off }D:
\]
the difference is $\lambda\log(1+|z|^2)$ in the contact disk and
$C_*-a\log(1+|z|^{-2})$ outside. Thus $E_C=\phi$ for $C\ge C_*$.
This proves $P_\theta[\psi]=\phi$ and $\phi\sim\psi$.

The singularity envelope depends only on bounded type: if
$\sigma\le\tau+M$, obstacle monotonicity compares its $C$ envelope with
the $(C+M)$ envelope of $\tau$. Increasing limits and the reverse comparison
give equality of the type envelopes. Hence $P_\theta[\phi]=\phi$;
the contact disk gives $\sup\phi=0$. For $\lambda/K<\epsilon<1$,
\[
 \epsilon\theta+\ddc\psi
 =(\epsilon K-\lambda)\omega_1+\epsilon b\omega_2+\lambda[D]>0.
\]
Since $\psi\le\phi$, this is an actual candidate below $\phi$ for the
smaller-background envelope, proving admissibility rather than assuming it.

The reference horizontal residual is $K\omega_1$ on $|z|^2<R$ and zero
outside. Matching derivatives yield no boundary atom. Consequently
\[
 \MA^{\mathrm{ext}}_\theta(\phi)
 =2Kb\,\mathbf1_{\{|z|^2<R\}}\,\mathrm{d}W.
\]
It has positive finite mass; the horizontal mass is
$K R/(1+R)=a$. Finally $\ddc\theta=(\ddc b)\wedge\omega_2\ne0$,
because a smooth harmonic function on a compact sphere is constant.
On a surface $\ddc(\theta+\ddc v)=\ddc\theta$ for every smooth $v$.
Thus none of its smooth positive perturbations is Gauduchon.

\section{Continuous residual and an actual weak product}

Let $d_{\mathrm{ch}}(x,z)$ be chordal distance on $\mathbb P^1$.
The kernel $G(x,z)=\log d_{\mathrm{ch}}(x,z)^2+1$ satisfies
\[
 \ddc_xG(x,z)=\delta_z-\omega_1,\qquad\int G(x,z)\omega_1(x)=0.
\]
The first identity is the one-variable logarithmic fundamental solution;
the second uses that squared chordal distance is uniform on $[0,1]$ under
normalized sphere area. For $f=A-a$, of mean zero, set
\[
 w(x)=\int_{\mathbb P^1}G(x,z)f(z)\omega_1(z).
\]
Uniform logarithmic integrability and $L^1$ continuity of translates imply
boundedness and continuity; explicitly
\[
 \|w\|_\infty\le\frac2e\|f\|_\infty,
 \qquad\int_0^1|\log t+1|\,\mathrm{d}t=\frac2e.
\]
Distributional differentiation gives $\ddc w=f\omega_1$ and mean zero.
A harmonic distribution on the sphere is constant, by inversion on the
nonconstant spherical harmonics, giving uniqueness of this normalization.

Equation \eqref{eq:pole} now gives
$\theta+\ddc u=\Omega+\lambda[D]>0$. The continuous $w$ and canonical
logarithm show that $u$ is genuine and $u\sim\phi$. Off $D$, its potential
and complex Hessian are continuous. Locally add a large positive quadratic
to make a bounded psh potential. Uniform convolution of this potential
and of its continuous Hessian shows, in the bounded recursion
\[
 (\ddc v)^2=\ddc(v\,\ddc v),
\]
that its product is the ordinary square of its continuous Hessian.
Polarization and expansion of the fixed smooth correction give $\Omega^2$.
The bounded recursion and its algebraic properties are
\cite[Section 2, (2.1), Theorem 2.1 and Corollary 2.2]{BT}.
The non-pluripolar definition removes $D$; continuous forms and their
squares give it zero mass. Extension by zero off $D$ therefore gives
$\Omega$ and $\Omega^2$ on all of $X$. No divisorial mixed product is retained.
This also applies to $u+t h$ whenever its residual remains positive.
The equation follows from \eqref{eq:density}; scalar uniqueness in the
verified admissible framework gives $c_\phi(F)=c$.

The weight in \eqref{eq:weight} is smooth. Thus its multiplication with
$\Omega$ is defined before any supporting probability is invoked.
The coefficient $A/b$ depends only on $x$, where $\omega_1$ has top degree;
distributional differentiation gives $\mathrm{d}((A/b)\omega_1)=0$.
This also follows by uniform smooth approximation of $A$. Hence
$\mathrm{d}Q=0$, and $\int\rho\Omega^2=\int(A/a)\omega_1=1$.
For small signed $t$, uniform positivity ensures $\Omega+t\ddc h>0$ and
\[
 \MA^{\mathrm{ext}}_\theta(u+t h)
 =\Omega^2+2t\Omega\wedge\ddc h+t^2(\ddc h)^2.
\]
This proves the asserted derivative and the continuous operator $L_u$.
Mixed components have zero wedge with the diagonal $\Omega$.
Distributional Stokes proves
\[
 \int L_uh\,\mathrm{d}\nu
 =2\int\ddc h\wedge Q=0.
 \]

\section{Full supporting probabilities and a signed-measure uniqueness proof}

Take any $\mu\in\suppset_\phi(F)$. For smooth $h$ and small signed $t$,
put
\[
 F_t=F+\log\frac{(\Omega+t\ddc h)^2}{\Omega^2}.
\]
This is continuous and $u+t h$ is an actual ordinary-type solution with
the same scalar $c$. Thus $B_\phi(F_t)=B_\phi(F)$. Uniform positivity gives
$F_t-F=tL_uh+O(t^2)$ in the uniform norm. Using \eqref{eq:support}
for both signs of $t$ proves
\begin{equation}\label{eq:annihilate}
 \int_XL_uh\,\mathrm{d}\mu=0\quad(h\in C^\infty(X,\mathbb R)).
\end{equation}
The perturbed $F_t$ need not satisfy \eqref{eq:density}; this step uses
the full supporting set, rather than a functional restricted to factored data.

\begin{lemma}\label{lem:resolvent}
Let $V$ be continuous and strictly positive on the round sphere, and $T$
a finite signed measure such that
\[
 \int(\Delta f-Vf)\,\mathrm{d}T=0\quad(f\in C^\infty(\mathbb P^1)).
\]
Then $T=0$.
\end{lemma}
\begin{proof}
For our normalized Fubini--Study form, $\Delta=\ddc/\omega$ is the
unit-sphere Laplacian, with degree-$\ell$ eigenvalue $-\ell(\ell+1)$.
The sphere heat semigroup $H_t$ multiplies the corresponding coefficients
by $e^{-t\ell(\ell+1)}$. Its series is smooth for $t>0$; the maximum
principle gives positivity and contraction for smooth initial functions.
Uniform density of smooth functions extends it to a strongly continuous
positive contraction on $C(\mathbb P^1)$.
Choose $M\ge\max V$. The resolvent
\[
 R_Mg=\int_0^\infty e^{-Mt}H_tg\,\mathrm{d}t
\]
has norm at most $1/M$ and $(M-\Delta)R_Mg=g$ in distributions.
For any continuous $g$, the equation $(\Delta-V)f=g$ becomes
\[
 f=R_M((M-V)f-g).
\]
This map contracts by $1-\min V/M<1$. Its fixed point is continuous
and has distributional Laplacian $\Delta f=Vf+g$ continuous.
The smooth $f_\varepsilon=H_\varepsilon f$ satisfy uniform convergence
of both $f_\varepsilon\to f$ and
$\Delta f_\varepsilon=H_\varepsilon(Vf+g)\to Vf+g$.
Hence $(\Delta-V)f_\varepsilon\to g$ uniformly. Testing against $T$
proves $\int g\,\mathrm{d}T=0$ for every continuous $g$, so $T=0$.
This argument imposes no prior density or Sobolev regularity on the measure.
\end{proof}

\begin{proposition}\label{prop:invariant}
The only Borel probability satisfying \eqref{eq:annihilate} is
$\nu=(A/a)\,\mathrm{d}W$.
\end{proposition}
\begin{proof}
Let $e_k$ be real spherical harmonics on the $y$ sphere, with $e_0=1$ and
$\Delta_2e_k=-\lambda_ke_k$. Only the constant mode has $\lambda_k=0$.
Their finite span is uniformly dense: restrictions of polynomials to the
sphere decompose into harmonics and are dense in its continuous functions.
Define finite signed measures on the first sphere by
\[
 \int f\,\mathrm{d}T_k
 =\int_X\frac{f(x)e_k(y)}{A(x)}\,\mathrm{d}\mu.
\]
Testing $L_u$ on $h=f(x)e_k(y)$ gives
\[
 \int\left(\Delta_1f-\lambda_k\frac A b f\right)\,\mathrm{d}T_k=0.
\]
For $k=0$, Laplacian inversion on smooth mean-zero sphere functions gives
$T_0=c_0\omega_1$. The $x$ marginal is $c_0A\omega_1$, and probability
normalization gives $c_0=1/a$. For every nonconstant mode, apply
Lemma~\ref{lem:resolvent} with the positive continuous $V=\lambda_k A/b$
to obtain $T_k=0$. Multiplication of test functions by $A$ then identifies
every mode of $\mu$: its zero mode is $(A/a)\omega_1$ and all other
modes vanish. Uniform harmonic density and density of finite product tests
give $\mu=(A/a)\omega_1\wedge\omega_2$. The preceding Stokes identity
verifies that this probability itself is invariant.
\end{proof}

Nonemptiness of the full supporting set, \eqref{eq:annihilate} and
Proposition~\ref{prop:invariant} prove \eqref{eq:nu}.
If $r\ge0$ is measurable with $\int r\Omega^2=1$, then $r$ is integrable
relative to $\mathrm{d}W$, because the continuous determinant is bounded
above and below. Its product with $\Omega$ is an order-zero current.
If $\ddc(r\Omega)=0$, then $r\Omega^2$ satisfies \eqref{eq:annihilate}.
It equals $\nu$, so $r=\rho$ almost everywhere. This proves normalized
weight uniqueness within the stated meaningful multiplication class.

\section{Approximation of the actual weighted currents}

For \eqref{eq:approx}, smooth residual adjoints exist and are unique
up to scale by Gauduchon's compact Hermitian conformal theorem
\cite[Section 3, Theorem 2 and its proof]{Gau}. Its hypotheses apply
to the connected compact surface and the positive smooth $\Omega_j$;
no even-Betti condition from the separate Theorem 1 is used.
Normalize as stated in Theorem~\ref{thm:main} and set
\[
 \nu_j=\rho_j\Omega_j^2,
 \qquad L_jh=\frac{2\Omega_j\wedge\ddc h}{\Omega_j^2}.
\]
These are probabilities and $\int L_jh\,\mathrm{d}\nu_j=0$.
Uniform tensor convergence to positive $\Omega$ gives uniform determinant
lower bounds and $L_jh\to L_uh$ uniformly for each fixed smooth $h$.
Every weak subsequential probability limit thus satisfies
\eqref{eq:annihilate}; Proposition~\ref{prop:invariant} forces it to be
$\nu$. Compactness of probabilities on $X$ proves convergence of the
whole sequence. Uniform determinants also give $F_j\to F$.

To prove weighted-current convergence, let $\xi$ be a smooth real
$(1,1)$ test form. The continuous functions
\[
 f_j=\frac{\Omega_j\wedge\xi}{\Omega_j^2}
 \longrightarrow f=\frac{\Omega\wedge\xi}{\Omega^2}
\]
converge uniformly. Therefore
\[
 \int\rho_j\Omega_j\wedge\xi
 =\int f_j\,\mathrm{d}\nu_j\longrightarrow
 \int f\,\mathrm{d}\nu=\int\rho\Omega\wedge\xi.
\]
This proves convergence of the actual weighted product, beyond convergence
of supporting probabilities. Likewise $q_j=\Omega_j^2/\mathrm{d}W$
converges uniformly to positive $q_0=\Omega^2/\mathrm{d}W$, so for every
continuous $g$,
\[
 \int g\rho_j\,\mathrm{d}W
 =\int(g/q_j)\,\mathrm{d}\nu_j\longrightarrow
 \int(g/q_0)\,\mathrm{d}\nu=\int g\rho\,\mathrm{d}W.
\]
Pointwise bounds for $\rho_j$ and convergence of its derivatives were
not required.

The class is nonempty. Apply the sphere heat semigroup to $A$ to obtain
smooth positive $A_j$ of mean $a$ converging uniformly to $A$.
Green inversion gives smooth mean-zero $w_j=w_j(x)$ with
$\ddc w_j=(A_j-a)\omega_1$. The kernel bound gives $w_j\to w$ uniformly,
and their residuals converge uniformly to $\Omega$.
Nonsymmetric paths are obtained by adding $\varepsilon_j h(x,y)$ with a
fixed smooth mean-zero genuinely mixed $h$ and $\varepsilon_j\to0$
sufficiently small to preserve positivity. The theorem proves independence
for every path satisfying \eqref{eq:approx}, not just these examples.
Each $u_j=\psi+w_j$ retains the pole and solves its smooth-density
equation with scalar $c$; it is not a globally smooth potential.

\section{A nonsmooth example and the pole-filling obstruction}

Take $a=1$, $\lambda=1/4$, $q=1$, and set
\[
 t=\frac{|Z_0|^2-|Z_1|^2}{|Z_0|^2+|Z_1|^2},
 \qquad b=2+t,\qquad H=A=\tfrac12+|t|,
\]
\[
 F=\log\bigl(2(\tfrac12+|t|)(2+t)\bigr).
\]
The variable $t$ is uniform on $[-1,1]$ under $\omega_1$. Thus
$\int A\omega_1=1$, $c=1$, $1/2\le A\le3/2$ and $1\le b\le3$.
The residual and density are Lipschitz and not $C^1$ at $t=0$;
the potential has a genuine unbounded logarithmic divisor.
The model contact radius squared is four, and
\[
 \tfrac12\theta+\ddc\psi
 =\tfrac38\omega_1+\tfrac b2\omega_2+\tfrac14[D]>0.
\]
Hence the example satisfies the exact admissibility antecedents. Its
weight and support are
\[
 \rho=\frac1{2(2+t)},\qquad\nu=(\tfrac12+|t|)\,\mathrm{d}W,
 \qquad\ddc\theta=-2t\,\mathrm{d}W\ne0.
\]
Direct integration gives
\[
 \int_X\Omega^2=4,\qquad
 \int_Xt\,\mathrm{d}\nu=0,\qquad
 \int_Xt\,\frac{\Omega^2}{4}=\frac5{24}.
\]
Thus the supporting probability is different from normalized
Monge--Amp\`ere volume.

In the same background instead take $H=A=a=1$, so $w=0$, and, for
$0<\delta\le1$, define
\[
 \psi_\delta=\lambda\log
 \frac{|Z_0|^2+\delta|Z_1|^2}{|Z_0|^2+|Z_1|^2},
 \qquad D_\delta=\omega_1+\ddc(\psi_\delta/\lambda).
\]
The potentials are globally smooth, decrease to $\psi$ as $\delta\downarrow0$,
converge in $L^1$ and smoothly off $D$. Direct differentiation gives
\[
 \frac{D_\delta}{\omega_1}
 =\frac{4\delta}{((1+\delta)+(1-\delta)t)^2}.
\]
This is positive of integral one and converges to $\delta_p$ at
$p=\{Z_0=0\}$. The positive raw solution forms are
\[
 \widehat\Omega_\delta=\theta+\ddc\psi_\delta
 =\omega_1+\lambda D_\delta+b\omega_2.
\]
Their normalized classical adjoints are
\[
 \widehat\rho_\delta=\frac1{2(1+\lambda)b},\qquad
 \widehat\nu_\delta
 =\frac{\omega_1+\lambda D_\delta}{1+\lambda}\wedge\omega_2.
\]
Their weighted forms are closed and their probability normalization is
immediate. Consequently
\[
 \widehat\nu_\delta\rightharpoonup
 \frac{\omega_1+\lambda\delta_p}{1+\lambda}\wedge\omega_2.
\]
For $\lambda=1/4$ this gives mass $1/5$ to $D$, whereas the genuine
full supporting probability is $\mathrm{d}W$. The weighted raw current
also retains the nonzero component
$\lambda[D]/(2(1+\lambda)b(p))$. It does not converge to $Q$.
This disproves unrestricted recovery from smooth pole-filling potential
approximations, even in this fixed non-Gauduchon background.
Those approximations fail the fixed-pole condition and uniform residual
convergence; their raw densities also fail uniform convergence to the target.
If pole smoothing is used, the prescribed residual construction removes
$\lambda D_\delta$ before computing the adjoint. The theorem only claims
independence for its explicitly specified nonempty approximation class.

\section{Scope and source applications}

The construction supplies a nonvacuous singular-reference sufficient-condition
answer with continuous data, a directly defined smooth weight on continuous
non-pluripolar residuals, a measure-valued linearization, a singleton full
supporting set and approximation-independent actual weighted currents.
The reference itself has residual degeneracy on an open set. The positive
target residual and uniform tensor approximation hypothesis are essential
to this proof; no general theorem for degenerate solution residuals or
arbitrary unbounded singularity types is asserted.

The indispensable external product input is the bounded local recursion in
\cite{BT}; the agreement with ordinary continuous-coefficient products is
proved above by uniform mollification and smooth-background expansion.
The classical residual adjoint input is the full proof of Theorem 2 in
\cite{Gau}, including its one-dimensional adjoint kernel and positivity.
Both sources use $\ddc=2i\partial\bar\partial$, four $\pi$ times our
operator. The recursion rescales homogeneously and the zero adjoint equation
is unchanged. No rough-weight multiplication theorem or general
continuous-coefficient invariant-measure uniqueness theorem is imported.
The original admissible normalization framework is explicitly identified
in Section 1, and every new use verifies the constructed model and type.

\begin{thebibliography}{9}
\bibitem{BT} E. Bedford and B. A. Taylor,
\emph{A new capacity for plurisubharmonic functions},
Acta Math. \textbf{149} (1982), 1--40.
\href{https://doi.org/10.1007/BF02392348}{doi:10.1007/BF02392348}.
\bibitem{Gau} P. Gauduchon,
\emph{Les m\'etriques standard d'une surface complexe compacte \`a premier
nombre de Betti pair}, Ast\'erisque \textbf{126} (1985), 129--135.
\href{https://www.numdam.org/article/AST_1985__126__129_0.pdf}{Author's original proof}.
\end{thebibliography}
\end{document}
